\chapter*{Introduction}
\addcontentsline{toc}{chapter}{\protect\numberline{}Introduction}

\section*{Motivation}
\addcontentsline{toc}{section}{\protect\numberline{}Motivation}

To protect business-critical applications from cybersecurity threats, security teams need to use a large set of application security tools: \gls{sast} scanners analyze source code for vulnerabilities, \gls{sca} tools detect issues in software dependencies, \gls{dast} solutions scan running applications for exploitable weaknesses, and secret scanning tools identify exposed credentials in code.

Further detection mechanisms need to be applied to software integration and deployment pipelines that are prone to their own classes of threats. Increasingly, security teams employ mechanisms that identify problems before they even reach the version control system, either in the form of IDE plugins or as security context tools integrated into AI coding assistants like Github Copilot or Claude Code.

In large enterprises, there may be dozens of such vulnerability detection tools, each producing findings through different \glspl{api}, data formats, and integration patterns \parencite{Anderson2020, Stallings2024}.

As application security architect in a large organization managing tens of thousands of code repositories and thousands of developers, I have witnessed firsthand how challenging it is to keep such a complex environment secure. Most notably:
\begin{itemize}
    \item Different detection tools of the same class capture different results as they employ different fundamental scanning paradigms
    \item The same vulnerability may be reported by multiple tools, requiring careful de-duplication
    \item Findings lack context regarding business criticality of the software project or its threat exposure, both of which are critical determinants of how urgently a vulnerability needs to be remediated
    \item TODO more
\end{itemize}

I argue that the challenge of consolidating application security data is, at its core, a data engineering problem. While the industry is addressing this with commercial products categorized as \gls{aspm} that offer dashboards and aggregated views, these solutions are proprietary, expensive, and inflexible. They introduce vendor lock-in and limit the ability of security teams to customize data pipelines, apply their own machine learning models, or integrate with internal systems on their own terms.

The most notable open-source alternative, DefectDojo, takes a different approach: it is a Django-based web application designed primarily as a vulnerability management interface. However, it was not built with a data-first architecture in mind. Its scaling capabilities are limited by the constraints of a traditional relational database backend, and its data ingestion model relies predominantly on push-based integrations or manual report uploads. Pull-based connectors that actively fetch data from security tool \glspl{api} are available only in the commercial DefectDojo Pro offering, further fragmenting the landscape between open and proprietary solutions.

What is missing is a vendor-agnostic, data-first framework that treats application security data as a first-class data engineering problem. One that provides a reusable architecture for ingestion, normalization, de-duplication, and serving of security findings at enterprise scale. This thesis aims to fill that gap.

\section*{Objective}
\addcontentsline{toc}{section}{\protect\numberline{}Objective}

The main goal of my thesis is to design and implement a data integration framework for enterprise application security. The framework consolidates security findings from heterogeneous tools into a unified data platform, enabling consistent normalization, de-duplication, triage, and correlation of vulnerabilities across business applications.

To achieve this goal, I define the following subgoals:

\begin{enumerate}
    \item \textbf{Literature review} (\autoref{ch:literature-review}): Survey existing literature across three domains: software development practices, application security, and data engineering. Identify the theoretical foundations and the gap that motivates this work.
    \item \textbf{Requirement specification} (\autoref{ch:requirement-analysis}): Analyze the application security domain to identify personas, data sources, required data transformations, data consumers, and non-functional requirements. Produce a requirements specification for the framework, with the full specification included in the appendix.
    \item \textbf{Architecture} (\autoref{ch:architecture}): Propose a reusable and extensible architecture that satisfies both analytical and transactional use cases.
    \item \textbf{Pipeline and model design} (\autoref{ch:design}): Design the data model and processing pipeline following the medallion architecture pattern~\cite{Strengholt2025}.
    \item \textbf{Implementation and validation} (\autoref{ch:implementation}): Use AI-assisted coding to build a reference implementation on the Databricks platform~\cite{Lee2024} and validate it against the requirements specification through automated testing with requirement traceability.
\end{enumerate}

\section*{Methodology}
\addcontentsline{toc}{section}{\protect\numberline{}Methodology}

This thesis follows a design-science research methodology as presented by~\textcite{Hevner2004, Peffers2007}, structured into three phases.

In the first phase, I conduct a \textbf{literature review} of academic and industry sources across three domains: software development, application security, and data engineering. This review establishes the theoretical foundation for the framework and includes vendor documentation for the technologies used in the implementation.

In the second phase, I perform a \textbf{requirements analysis and design}. I analyze the application security landscape to identify the key personas, data sources, required transformations, and non-functional requirements. Requirements are prioritized using the MoSCoW method. Based on the requirements, I propose an architecture and design a data model with a medallion-based pipeline.

In the third phase, I build a \textbf{reference implementation} on the Databricks platform to demonstrate feasibility of the proposed design. I validate the implementation through automated tests with requirement traceability, mapping individual tests to requirements, and through data quality expectations enforced by \gls{dltables}.

The thesis delivers three distinct contributions:
\begin{enumerate}
    \item A \textbf{requirements specification} for an application security data platform that other organizations can adapt to their own needs.
    \item A \textbf{reusable design} covering architecture, data model, and pipeline that is not tied to a specific platform or vendor.
    \item A \textbf{reference implementation} on Databricks that proves the feasibility of the design and serves as a starting point for production deployments.
\end{enumerate}

\section*{Structure}
\addcontentsline{toc}{section}{\protect\numberline{}Structure}

\autoref{ch:literature-review} of the thesis researches relevant literature on software development, application security, and data engineering, and identifies the research gap. \autoref{ch:requirement-analysis} defines the personas, data sources, transformations, consumers, and non-fun\-ctio\-nal requirements, organized in a coherent requirements specification. \autoref{ch:architecture} presents the platform architecture, including technology selection and component design. \autoref{ch:design} details the medallion-based data model and pipeline transformations across the bronze, silver, and gold layers. \autoref{ch:implementation} describes the reference implementation on Da\-ta\-bricks, covering connector development, pipeline orchestration, and testing. The Conclusion evaluates the outcomes against the objectives defined above, discusses limitations, and suggests directions for future work.